% Outline: Background (1.5 page cap, >=2 figures; defines the paper's vocabulary)
\section{Background}

\subsection{Serving, Scheduling, and Head-of-Line Blocking}
Modern LLM engines admit requests continuously and batch them at
token granularity; continuous batching~\cite{orca2022} and vLLM's paged
KV-cache~\cite{vllm2023} are the substrate for this work. The scheduler decides which waiting request joins
the running batch next. First-come-first-served is the default, and its
weakness is classical: a long-running request admitted early delays every
short request behind it. A 2026 survey of latency in LLM agent systems
independently lists FIFO-induced head-of-line blocking among live system
bottlenecks for agentic traffic~\cite{park2026minimizing}, while placing
learned scheduling in its future-work section --- the gap this paper works in.

\subsection{Learned Output-Length Prediction}
Shortest-job-style ordering requires knowing output lengths in
advance. Since generation length is stochastic, a body of work predicts it:
S\textsuperscript{3} classifies responses into length buckets; the
learning-to-rank line~\cite{fuLTR2024} showed that \emph{relative order},
not absolute length, is what scheduling needs, training a small proxy with a
ranking loss; PARS~\cite{tao2026ranking} refined this with pairwise margin
training over a BERT encoder. All of these read the prompt (or scalar request
statistics); none consume tool-schema text.

\subsection{System Substrate: Gateway and Decision Service}
Our serving stack (Fig.~\ref{fig:arch}) fronts the engine with an
LLM gateway (VeloxMesh) that owns routing, admission, and a scheduler hook
with a hard design contract: scoring must return within 15\,ms, after which
the gateway \emph{fails open} to arrival order. The gateway consults a
separate decision service over HTTP (Fig.~\ref{fig:components}); the service
hosts the Ranker and applies the Reliability Gate on the gateway's behalf.
This fail-open contract is load-bearing for everything that follows: an
unreliable or slow prediction never blocks the request path.

\begin{figure*}[t]
  \centering
  \begin{tikzpicture}[node distance=6.5mm and 9mm]
    \node[environment component, text width=19mm] (client)
      {Agent clients\\[-1pt]{\sub{OpenCode}}};
    \node[environment component, right=of client, text width=21mm] (gw)
      {Gateway\\[-1pt]{\sub{admission, routing}}};
    \node[proposed component, right=of gw, text width=25mm] (dec)
      {Decision service\\[-1pt]{\sub{Ranker + Reliability Gate}}};
    \node[environment component, right=15mm of dec, text width=19mm] (fb)
      {Fallback\\[-1pt]{\sub{arrival order}}};
    \node[proposed component, below=11mm of dec, text width=25mm] (order)
      {Queue order\\[-1pt]{\sub{policy under test}}};
    \node[environment component, below=11mm of gw, text width=21mm] (engine)
      {vLLM engine\\[-1pt]{\sub{Qwen3.5-9B}}};

    \draw[data flow] (client) -- node[numbered step, pos=.5] {1} (gw);
    \draw[control flow] (gw) -- node[numbered step, pos=.5] {2} (dec);
    \draw[control flow] (dec) -- node[numbered step, pos=.5] {3} (order);
    \draw[control flow] (order) -- node[numbered step, pos=.5] {4} (engine);
    \draw[data flow] (gw) -- node[numbered step, pos=.5] {5} (engine);
    \draw[degraded flow] (dec) -- node[warn label, pos=.5, above=1pt]
      {$>15$\,ms} (fb);

    \begin{scope}[on background layer]
      \node[system boundary, fit=(gw)(dec)(order)(engine)] (sys) {};
    \end{scope}
    \node[boundary title, anchor=north west]
      at ($(sys.north west)+(3pt,-2pt)$) {serving stack under study};
  \end{tikzpicture}
  \caption{System architecture. Solid arrows carry requests, dashed arrows
  carry scheduling decisions, and the orange path is the gateway's fail-open
  behavior when scoring exceeds its 15\,ms budget. Shaded components are
  contributed by this work; the rest is the environment they run in.}
  \label{fig:arch}
\end{figure*}

\begin{figure*}[t]
  \centering
  \begin{tikzpicture}[node distance=7mm and 12mm]
    \node[environment component, text width=19mm] (req)
      {Request\\[-1pt]{\sub{prompt + schema}}};
    \node[proposed component, right=of req, text width=24mm] (cls)
      {Stratum classifier\\[-1pt]{\sub{hash vs vocabulary}}};
    \node[proposed component, above right=3.5mm and 22mm of cls, text width=26mm]
      (rank) {Ranker\\[-1pt]{\sub{BERT fp16, batch $\leq$ 8}}};
    \node[proposed component, right=of rank, text width=21mm] (ord)
      {Rank order\\[-1pt]{\sub{trusted slots only}}};
    \node[environment component, below right=3.5mm and 22mm of cls,
      text width=26mm] (queue) {Fallback queue\\[-1pt]{\sub{arrival order}}};
    \node[environment component, right=of queue, text width=21mm] (served)
      {Served\\[-1pt]{\sub{slot preserved}}};

    % The branch is routed orthogonally from a single split point rather than
    % as two diagonals. One decision, two outcomes -- the geometry should say
    % that, and right angles keep the edge labels off the boxes.
    \coordinate (split) at ($(cls.east)+(9mm,0)$);
    \draw[draw=black!80, line width=.8pt] (cls.east) -- (split);
    \draw[control flow] (split) |- (rank.west);
    \draw[degraded flow] (split) |- (queue.west);
    \draw[data flow] (req) -- (cls);
    \draw[data flow] (rank) -- (ord);
    \draw[data flow] (queue) -- (served);

    \node[fact label, above right=0.5mm and 1mm of split] {S3 / S4};
    \node[warn label, below right=0.5mm and 1mm of split, align=left]
      {S1 / S2 / zero-tool\\no Ranker invocation};
  \end{tikzpicture}
  \caption{Decision path. The stratum classifier is a vocabulary hash, so
  requests the gate cannot vouch for are answered without invoking the model at
  all. Abstained requests keep their arrival slots; trusted requests are
  reordered only among themselves.}
  \label{fig:components}
\end{figure*}

\subsection{Vocabulary}
Throughout, the \emph{Ranker} is the fine-tuned BERT model scoring a
request's expected output length; the \emph{Reliability Gate} is the runtime
decision of whether that score may influence queue order; \emph{Fallback}
means serving in arrival order because the gate withheld trust.
\emph{Ranking Tau} is Kendall's $\tau_b$~\cite{kendall1938} between predicted and true
output-length order on held-out data. \emph{Cold-Start Transfer} measures
Ranking Tau on test strata whose tool sets were unseen in training, within
the same workload family: S1 (seen combination), S2 (new combination of seen
tools), S3 (partially new tools), S4 (all-new tools). It is distinct from
\emph{Cross-Workload Transfer}, which changes the workload family entirely.
